\documentclass[aspectratio=169,11pt]{beamer}
\usetheme{Madrid}
\usecolortheme{seahorse}
\usefonttheme{professionalfonts}
\usepackage[T1]{fontenc}
\usepackage[utf8]{inputenc}
\usepackage[spanish,es-tabla]{babel}
\usepackage{lmodern}
\usepackage{microtype}
\usepackage{amsmath}
\usepackage{hyperref}

\hypersetup{colorlinks=true,urlcolor=blue}
\setbeamertemplate{navigation symbols}{}
\setbeamertemplate{footline}{}

\title[Premio Nobel de Física 2023]{Premio Nobel de Física 2023:\\pulsos de attosegundos y dinámica electrónica}
\subtitle{Agostini, Krausz y L'Huillier\\Métodos para generar pulsos de attosegundos y estudiar electrones en la materia [1]}
\author{Presentación de investigación}
\date{21 de junio de 2026}

\begin{document}

\begin{frame}
\titlepage
\end{frame}

\begin{frame}{Antecedentes: de la femtoquímica a la attociencia}
\begin{itemize}
\item La femtoquímica resolvió movimientos nucleares y estados de transición; la escala electrónica exige $10^{-18}$ s [1,2].
\item Un attosegundo es $10^{-18}$ s; la unidad atómica de tiempo es $24.2$ as, comparable con correlaciones electrónicas [1].
\item La condición física para un pulso corto es una banda espectral ancha y coherente: $\Delta t\,\Delta E \sim \hbar$.
\item La generación de armónicos altos (HHG) en gases nobles dio el salto: L'Huillier observó en 1987 una meseta de armónicos impares bajo láser IR intenso [3].
\item Esa meseta aportó frecuencias \textsc{xuv} suficientes para sintetizar trenes de attopulsos.
\end{itemize}
\end{frame}

\begin{frame}{Mecanismo físico: HHG y modelo de tres pasos}
\begin{columns}[T,totalwidth=\textwidth]
\begin{column}{0.56\textwidth}
\begin{enumerate}
\item Ionización por túnel: el campo láser deforma el potencial atómico.
\item Aceleración: el electrón libre gana energía en el campo oscilante.
\item Recolisión y recombinación: el electrón regresa al ion y emite un fotón \textsc{xuv}.
\end{enumerate}
\vspace{0.05in}
\[
E_c \approx I_p + 3.17U_p
\]
\small La energía de corte depende del potencial de ionización y del potencial ponderomotriz [4--6].
\end{column}
\begin{column}{0.40\textwidth}
\small
\textbf{Claves conceptuales}
\begin{itemize}
\item La HHG es coherente y sincronizada con el láser impulsor.
\item El plateau permite duración temporal extrema.
\item La fase de los armónicos importa tanto como su intensidad.
\item La teoría de Lewenstein-L'Huillier-Corkum conectó el modelo semiclasico con una descripción cuántica [6].
\end{itemize}
\end{column}
\end{columns}
\end{frame}

\begin{frame}{Desarrollo experimental y resultados premiados}
\begin{itemize}
\item Agostini: experiencia en ionización por encima del umbral y metrología interferométrica; su grupo generó trenes de pulsos de 250 as y los caracterizó con RABBITT [7,8].
\item RABBITT mide fases relativas mediante interferencia entre rutas de dos fotones: armónico + absorción IR frente a armónico vecino + emisión IR [8].
\item Krausz: pulsos de pocos ciclos, selección de armónicos de corte y streaking; su grupo aisló un pulso único de 650 as [9,10].
\item Streaking: un pulso \textsc{xuv} libera electrones y un campo IR sincronizado codifica el instante de emisión en la energía final [10].
\item L'Huillier: descubrió la meseta de HHG, desarrolló la teoría y consolidó el uso de trenes de attopulsos para medir retardos de fotoionización [3,6,15].
\end{itemize}
\end{frame}

\begin{frame}{Dinámica electrónica: resultados científicos}
\begin{itemize}
\item El efecto fotoeléctrico dejó de tratarse como instantáneo: en neón se midió un retardo de 21 as entre electrones 2s y 2p [11].
\item La discrepancia con teoría impulsó experimentos de mayor resolución espectral; el grupo de L'Huillier separó canales de shake-up y obtuvo acuerdo con teoría de muchos cuerpos [12--15].
\item En agua líquida se midieron retardos de 50 a 70 as respecto al vapor, atribuidos a la solvatación y al paisaje potencial del líquido [16].
\item En tungsteno se observaron retardos de alrededor de 100 as entre estados 4f localizados y estados de banda de conducción [17].
\item El observable central es fase-tiempo-energía: revela correlación electrónica, apantallamiento, transporte inicial y coherencia.
\end{itemize}
\end{frame}

\begin{frame}{Implicaciones y avances posteriores}
\begin{itemize}
\item Materiales: HHG en sólidos, corrientes petahertz, excitones, topología, fase de Berry y emisión dependiente de bandas [20,21].
\item Química: migración de carga y acoplamientos no adiabáticos antes del movimiento nuclear; base de la ``attoquímica'' [23].
\item Fuentes: pulsos de rayos X blandos hacia la ventana de agua y reportes recientes de 19.2 as por debajo de la unidad atómica de tiempo [18].
\item Metrología: absorción transitoria completamente attosegundo y fuentes \textit{table-top} que podrían replicarse en más laboratorios [19].
\item Aplicaciones adyacentes: huellas moleculares infrarrojas para muestras biomédicas y detección de patrones complejos, todavía con retos de validación clínica [22].
\end{itemize}
\end{frame}

\begin{frame}{Referencias}
\tiny
\begin{enumerate}
\item Royal Swedish Academy of Sciences. \textit{Scientific Background}. NobelPrize.org, 2023. \url{https://www.nobelprize.org/uploads/2023/10/advanced-physicsprize2023-2.pdf}
\item Nobel Prize Outreach. \textit{Press release: The Nobel Prize in Physics 2023}. \url{https://www.nobelprize.org/prizes/physics/2023/press-release/}
\item M. Ferray et al., \textit{J. Phys. B} 21, L31, 1988.
\item A. L'Huillier, K. J. Schafer y K. C. Kulander, \textit{J. Phys. B} 24, 3315, 1991.
\item P. B. Corkum, \textit{Phys. Rev. Lett.} 71, 1994, 1993.
\item M. Lewenstein et al., \textit{Phys. Rev. A} 49, 2117, 1994.
\item P. Agostini et al., \textit{Phys. Rev. Lett.} 42, 1127, 1979.
\item P. M. Paul et al., \textit{Science} 292, 1689, 2001.
\item M. Nisoli et al., \textit{Opt. Lett.} 22, 522, 1997.
\item M. Hentschel et al., \textit{Nature} 414, 509, 2001.
\item M. Schultze et al., \textit{Science} 328, 1658, 2010.
\item L. R. Moore et al., \textit{Phys. Rev. A} 84, 061404(R), 2011.
\item J. M. Dahlström et al., \textit{Phys. Rev. A} 86, 061402(R), 2012.
\item J. Feist et al., \textit{Phys. Rev. A} 89, 033417, 2014.
\item M. Isinger et al., \textit{Science} 358, 893, 2017.
\item I. Jordan et al., \textit{Science} 369, 974, 2020.
\item A. L. Cavalieri et al., \textit{Nature} 449, 1029, 2007.
\item F. Ardana-Lamas et al., arXiv:2510.04086, 2025.
\item M. Volkov et al., arXiv:2512.09585, 2025.
\item M. F. Ciappina, arXiv:2510.15207, 2025.
\item S. Martín-Domene et al., arXiv:2606.03945, 2026.
\item M. Zigman et al., \textit{Annals of Oncology} 33, S580, 2022.
\item F. Krausz y M. Ivanov, \textit{Rev. Mod. Phys.} 81, 163, 2009.
\end{enumerate}
\end{frame}

\end{document}
